\scanpage{161}
Si on a sur $\underline U$ et $\underline V$ des augmentations continues,
les hom. compatibles avec l'augmentation
correspondent aux foncteurs permis qui transforment
le $\underline U$ module trivial \deleted{\unsure{aug}} $k$ en le $\underline V$-module trivial $\underline k$.
Transitivité, etc.

\deleted{Supposons donné un foncteur}

\uline{Autres structures supplémentaires}

\inserted{Soit $k$ un \deleted{\unsure{fonct}}\inserted{corps}.} Supposons donné un foncteur $F^\vee\ \unsure{\mathcal K|\mathscr C|\mathcal C}^{\underline U}\longrightarrow\mathcal K^{\underline U}$
\deleted{tel que l'on ait une \unsure{identité}}
associant à \deleted{tout} $V\in K^{\underline U}$ une structure
de $\underline U$-module permise sur le dual $V'$ de $V$,
\deleted{De telle façon que}
i.e. un foncteur tel que
\diagram{\begin{tikzcd}[ampersand replacement=\&] \underline K^{\underline U}\arrow[r,"F^\vee"]\arrow[d] \& \underline K^{\underline U}\arrow[d] \& {} \\ \underline K^{k}\arrow[r,"F^\vee"{name=lowerF}] \& \underline K^{k} \& \breve F(W)=V'\arrow[from=lowerF,to=2-3,bend right=15] \end{tikzcd}}
soit commutatif. \deleted{\unsure{Cela comme}}
\deleted{$V'$ est un $\underline U^\circ$-module ($\underline U^\circ$ inverse de $\underline U$) cela signifie}
On veut que de plus,
$F^\vee F^\vee={}\underline{\text{identité}}$. \deleted{Ceci définit, par} \deleted{Il revient au $=$ de}
\deleted{\unsure{se donner}, sur tout} Soit alors $V$ un
\deleted{tout} module permis sur $\underline U^\circ$
(\unsure{l'opposé|l'opposée|l'inverse} de $\underline U$) $V'$ est un module sur
$\underline U$ (transposition) donc $(V')'$ est un
module sur $\underline U$. On a ainsi un
\deleted{\unsure{contrav.}} foncteur permis des $\underline U^\circ$-

% archive sheet number: 160
